\clearpage\pagestyle{otMain}\fontsize{11}{15}\selectfont\corrigesChapitre{14}{Trigonométrie du triangle rectangle}
\begin{blocLivre}{cor-14-1}\begin{corrige}{1}{Nommer les côtés}\demarche{Commencer par l’hypoténuse.}
Le côté opposé est $\res{AC}$, l’adjacent est $\res{AB}$ et l’hypoténuse est $\res{BC}$. AB touche B mais n’est pas l’hypoténuse.
\label{fig-cor-14-1}\figureGeo[corrige]{trigo-01}
\end{corrige}\end{blocLivre}
\begin{blocLivre}{cor-14-2}\begin{corrige}{2}{Choisir le cosinus}\demarche{Écrire les noms des côtés avant les nombres.}
$\cos B=AB/BC=6/10=\res{0{,}6}$. Le rapport utilise l’adjacent et l’hypoténuse.
\label{fig-cor-14-2}\figureGeo[corrige]{trigo-02}
\end{corrige}\end{blocLivre}
\begin{blocLivre}{cor-14-3}\begin{corrige}{3}{Choisir le sinus}\demarche{Repérer le côté qui ne touche pas l’angle B.}
$\sin B=AC/BC=8/10=\res{0{,}8}$. AC est opposé à B.
\label{fig-cor-14-3}\figureGeo[corrige]{trigo-03}
\end{corrige}\end{blocLivre}
\begin{blocLivre}{cor-14-4}\begin{corrige}{4}{Choisir la tangente}\demarche{La tangente ne fait pas intervenir l’hypoténuse.}
$\tan B=AC/AB=8/6=\res{4/3}$. Une tangente peut dépasser 1, contrairement au sinus ou au cosinus d’un angle aigu.
\label{fig-cor-14-4}\figureGeo[corrige]{trigo-04}
\end{corrige}\end{blocLivre}
\begin{blocLivre}{cor-14-5}\begin{corrige}{5}{Une première longueur}\demarche{Isoler la longueur dans le rapport.}
$AC=BC\sin30^\circ=12\times0{,}5=\res{6\text{ cm}}$. Le côté opposé est la moitié de l’hypoténuse pour cet angle.
\label{fig-cor-14-5}\figureGeo[corrige]{trigo-05}
\end{corrige}\end{blocLivre}
\begin{blocLivre}{cor-14-6}\begin{corrige}{6}{Un côté adjacent}\demarche{Relier la précision à l’unité de travail.}
$AB=9\cos40^\circ\approx6{,}8944$ cm, donc $\res{AB\approx6{,}9\text{ cm}}$ au millimètre. Le millimètre représente 0,1 cm.
\label{fig-cor-14-6}\figureGeo[corrige]{trigo-06}
\end{corrige}\end{blocLivre}
\begin{blocLivre}{cor-14-7}\begin{corrige}{7}{Une hypoténuse}\demarche{Résoudre l’égalité avant de calculer.}
$\cos50^\circ=7/BC$, donc $BC=7/\cos50^\circ\approx\res{10{,}9\text{ cm}}$. Multiplier par le cosinus aurait donné un côté trop court.
\label{fig-cor-14-7}\figureGeo[corrige]{trigo-07}
\end{corrige}\end{blocLivre}
\begin{blocLivre}{cor-14-8}\begin{corrige}{8}{Un angle par cosinus}\demarche{Utiliser la fonction inverse en mode degrés.}
$\cos B=0{,}6$. La fonction inverse donne $B\approx53{,}13^\circ$, donc $\res{B\approx53^\circ}$. L’angle C vaut alors environ 37 degrés.
\label{fig-cor-14-8}\figureGeo[corrige]{trigo-08}
\end{corrige}\end{blocLivre}
\begin{blocLivre}{cor-14-9}\begin{corrige}{9}{Un angle par tangente}\demarche{Choisir le rapport contenant les deux données.}
$\tan B=5/12$, donc $B=\arctan(5/12)\approx\res{22{,}6^\circ}$. L’opposé est inférieur à l’adjacent : un angle inférieur à 45 degrés est cohérent.
\label{fig-cor-14-9}\figureGeo[corrige]{trigo-09}
\end{corrige}\end{blocLivre}
\begin{blocLivre}{cor-14-10}\begin{corrige}{10}{Changer d’angle}\demarche{Réidentifier les côtés quand l’angle change.}
$\sin B=12/15=4/5$ et $\cos C=12/15=4/5$. Le côté opposé à B est adjacent à C ; les deux angles sont complémentaires.
\label{fig-cor-14-10}\figureGeo[corrige]{trigo-10}
\end{corrige}\end{blocLivre}
\begin{blocLivre}{cor-14-11}\begin{corrige}{11}{La hauteur d’un bâtiment}\demarche{Séparer hauteur au-dessus de la visée et hauteur totale.}
Au-dessus de l’appareil, la hauteur vaut $25\tan32^\circ\approx15{,}6217$ m. En ajoutant 1,6 m, la hauteur totale est $\res{17{,}2\text{ m}}$ au décimètre. Oublier la hauteur de visée sous-estimerait le résultat.
\label{fig-cor-14-11}\figureGeo[corrige]{trigo-batiment}
\end{corrige}\end{blocLivre}
\begin{blocLivre}{cor-14-12}\begin{corrige}{12}{Une rampe et son angle}\demarche{Distinguer pourcentage de pente et mesure angulaire.}
$B=\arctan(0{,}8/10)\approx\res{4{,}6^\circ}$. La pente est $\res{8\,\%}$. 8 \% ne signifie donc pas un angle de 8 degrés.
\label{fig-cor-14-12}\figureGeo[corrige]{trigo-12}
\end{corrige}\end{blocLivre}
\begin{blocLivre}{cor-14-13}\begin{corrige}{13}{Un câble}\demarche{Calculer les deux projections de l’hypoténuse.}
$AC=18\sin65^\circ\approx\res{16{,}3\text{ m}}$ et $AB=18\cos65^\circ\approx\res{7{,}6\text{ m}}$. Les calculs non arrondis vérifient Pythagore.
\label{fig-cor-14-13}\figureGeo[corrige]{trigo-13}
\end{corrige}\end{blocLivre}
\begin{blocLivre}{cor-14-14}\begin{corrige}{14}{Un toit symétrique}\demarche{Exploiter la symétrie pour obtenir un triangle rectangle.}
La hauteur partage la base en deux segments de 4 m. Elle vaut $4\tan35^\circ\approx\res{2{,}80\text{ m}}$. Un versant mesure $4/\cos35^\circ\approx\res{4{,}88\text{ m}}$. On utilise 4 m, pas la base entière.
\label{fig-cor-14-14}\figureGeo[corrige]{trigo-toit}
\end{corrige}\end{blocLivre}
\begin{blocLivre}{cor-14-15}\begin{corrige}{15}{Diagnostiquer une erreur}\demarche{Contrôler les côtés du rapport, pas seulement le résultat.}
AB est adjacent et AC opposé à B ; leur rapport est la tangente. $\tan38^\circ=AC/7$, donc $\res{AC=7\tan38^\circ\approx5{,}5\text{ cm}}$. Le sinus exigerait l’hypoténuse au dénominateur.
\label{fig-cor-14-15}\figureGeo[corrige]{trigo-15}
\end{corrige}\end{blocLivre}
